\section*{Conclusion et recommandations}
\addcontentsline{toc}{section}{Conclusion et recommandations}
\markboth{Conclusion et recommandations}{Conclusion et recommandations}

Ce mémoire avait pour objectif d'évaluer l'impact d'une exposition durable
aux transferts de migrants sur la pauvreté multidimensionnelle des enfants
au Sénégal, en exploitant les deux éditions de l'Enquête harmonisée sur
les conditions de vie des ménages. La stratégie d'identification repose
sur l'appariement par score de propension~: l'effet est estimé sur le
nombre de privations de la seconde édition, séparément dans chaque groupe
d'âge, le score incluant le niveau initial de privations de chaque enfant
et les caractéristiques prédéterminées du ménage.

Le résultat principal est un effet protecteur modeste et localisé. Chez
les enfants de 5 à 14 ans, une exposition durable aux transferts réduit le
nombre de dimensions en privation de 0,19, effet significatif au seuil de
10\,\%~; l'effet est de même signe mais imprécis chez les 0 à 4 ans, et
l'effectif des adolescents suivis ne permet aucune conclusion. Cette
réduction est cohérente avec le canal budgétaire des envois, l'écart entre
bénéficiaires et non-bénéficiaires étant le plus marqué sur la sécurité
alimentaire, et elle bute sur les privations d'infrastructure, que le
pouvoir d'achat ne lève pas.

Au regard des hypothèses de départ, H1 est partiellement validée~: une
exposition durable aux transferts réduit la pauvreté multidimensionnelle
des enfants d'âge scolaire, sans que l'effet soit établi aux autres âges.
H2 l'est également en partie~: l'effet se concentre sur les privations qui
répondent au budget courant, la sécurité alimentaire en tête, et croît
fortement avec le montant reçu, seuls les transferts substantiels
réduisant les privations.

\subsection*{Implications de politique économique}

Les résultats obtenus conduisent aux recommandations suivantes, à
l'intention des décideurs politiques. Chacune part d'un fait établi par
l'analyse et en tire une action, de la consolidation du canal par lequel
l'effet protecteur passe jusqu'à la levée des contraintes qui l'empêchent
de s'étendre.

\begin{enumerate}
  \item Consolider le canal alimentaire et l'étendre aux ménages qui
        n'en bénéficient pas. L'insécurité alimentaire est la seule
        dimension où l'effet des transferts est établi, et l'écart entre
        bénéficiaires et non-bénéficiaires y est le plus large de tous
        les indicateurs~: les envois protègent d'abord l'assiette des
        enfants. Les filets sociaux, les cantines scolaires et le suivi
        nutritionnel communautaire devraient donc cibler en priorité les
        enfants de ménages non bénéficiaires des communes d'émigration,
        privés de cet amortisseur privé.

  \item Cibler les enfants de 0 à 4 ans par l'offre de services plutôt
        que par le revenu. Chez les plus jeunes, l'effet protecteur ne se
        matérialise pas sur le nombre total de privations, leurs manques
        dominants, l'acte de naissance et l'environnement sanitaire du
        logement, passant par des démarches et des équipements que
        l'argent reçu ne déclenche pas. La consultation prénatale et
        postnatale, l'appui à l'état civil, la vaccination et les
        structures de garde méritent un renforcement ciblé dans les
        communes d'émigration.

  \item Maintenir les ménages bénéficiaires dans le champ de la
        protection sociale. L'effet protecteur mesuré est modeste, de
        l'ordre d'un cinquième de privation, et l'analyse par montant
        montre que seuls les transferts de la moitié supérieure de la
        distribution, au-delà d'environ 1,25 million de FCFA par an,
        réduisent effectivement les privations. Recevoir de l'étranger ne
        saurait donc valoir exclusion des dispositifs publics, et les
        ménages aux envois modestes doivent être traités, du point de vue
        du ciblage, comme des ménages non couverts.

  \item Augmenter les montants effectivement reçus en réduisant les
        coûts de transaction. Puisque l'effet croît avec le montant, tout
        franc perdu en frais d'envoi est une protection perdue pour
        l'enfant. La concurrence régulatoire sur les frais des canaux
        formels et l'expansion du transfert par téléphonie mobile
        augmenteraient les sommes disponibles et, avec elles, la capacité
        des envois à desserrer le budget courant.

  \item Compléter les transferts par l'offre publique de services de
        base. L'eau, l'assainissement et la santé relèvent
        d'investissements collectifs que le pouvoir d'achat privé ne
        remplace pas, comme le montre la privation de santé, quasi
        universelle dans les deux groupes. Les programmes d'accès
        devraient viser prioritairement les communes rurales
        d'émigration, en tirant parti du réseau associatif des diasporas
        pour amplifier l'effet sur les dimensions les plus rigides.
\end{enumerate}
